\documentclass[a4paper]{arrowhead}

\usepackage[yyyymmdd]{datetime}
\usepackage{etoolbox}
\usepackage[utf8]{inputenc}
\usepackage[T1]{fontenc}
\usepackage{multirow}
\usepackage{longtable}
\usepackage{tikz}
\usepackage{adjustbox}
\AtBeginEnvironment{longtable}{\small}
\BeforeBeginEnvironment{tikzpicture}{\begin{adjustbox}{max width=\textwidth}}
\AfterEndEnvironment{tikzpicture}{\end{adjustbox}}
\usetikzlibrary{positioning,shapes.multipart,arrows.meta}

\renewcommand{\dateseparator}{-}
\setlength{\emergencystretch}{3em}
\lstset{language=HTTP,breakatwhitespace=false,columns=fullflexible}

%% Special references
\newcommand{\mref}[2][]{{\textcolor{ArrowheadPurple}{\hyperref[sec:model:#2]{\ifstrempty{#1}{#2}{#1}}}}}
\newcommand{\oref}[1]{{\textcolor{ArrowheadBlue}{\hyperref[sec:operations:#1]{#1}}}}
\newcommand{\pdef}[1]{{\textcolor{ArrowheadGrey}{#1\label{sec:model:primitives:#1}\label{sec:model:primitives:#1s}\label{sec:model:primitives:#1es}}}}
\newcommand{\pref}[1]{{\textcolor{ArrowheadGrey}{\hyperref[sec:model:primitives:#1]{#1}}}}

% as template_sd.tex; the input and output arguments are passed already formatted (\mref or \pref)
\newrobustcmd\osubsection[3]{
  \addtocounter{subsection}{1}
  \addcontentsline{toc}{subsection}{\protect\numberline{\thesubsection}operation \textcolor{ArrowheadBlue}{#1}}
  \subsection*{
    \thesubsection\quad
    operation
    \textcolor{ArrowheadBlue}{#1}
    (\notblank{#2}{#2}{})
    \notblank{#3}{: #3}{}
  }
  \label{sec:operations:#1}
}
\newrobustcmd\msubsection[2]{
  \addtocounter{subsection}{1}
  \addcontentsline{toc}{subsection}{\protect\numberline{\thesubsection}#1 \textcolor{ArrowheadPurple}{#2}}
  \subsection*{\thesubsection\quad#1 \textcolor{ArrowheadPurple}{#2}}
  \label{sec:model:#2} \label{sec:model:#2es}
}
%%

\begin{document}
\sloppy

%% Arrowhead Document Properties
\ArrowheadTitle{authorizationToken}
\ArrowheadServiceID{authorizationToken}
\ArrowheadType{Service Description}
\ArrowheadTypeShort{SD}
\ArrowheadVersion{5.2.2}
\ArrowheadDate{2026-10-22}
\ArrowheadAuthor{Jerker Delsing}
\ArrowheadStatus{RELEASE CANDIDATE}
\ArrowheadContact{jerker.delsing@ltu.se}
\ArrowheadFooter{\href{www.arrowhead.eu}{www.arrowhead.eu}}
\ArrowheadSetup
%%

%% Front Page
\begin{center}
  \vspace*{1cm}
  \huge{\arrowtitle}

  \vspace*{0.2cm}
  \LARGE{\arrowtype}
  \vspace*{1cm}

  \Large{Service ID: \textit{"\arrowid"}}

  \vspace*{\fill}

  % Front Page Image

  \vspace*{1cm}
  \vspace*{\fill}

  % Front Page Abstract
  \begin{abstract}
    This document describes the \texttt{authorizationToken} service, version 1.0.0, produced by the ConsumerAuthorization core system of Arrowhead 5.2.2.
  \end{abstract}

  \vspace*{1cm}
\end{center}

\newpage
%%

%% Table of Contents
\tableofcontents
\newpage
%%

\section{Overview}
\label{sec:overview}

This document describes the \texttt{authorizationToken} service, which provides authorization tokens: a consumer obtains a token for a specific provider, service and operation; the provider verifies it.
It is produced by the ConsumerAuthorization system \cite{ca-sysd} in service version 1.0.0.

The rest of this document is organized as follows.
In Section \ref{sec:operations}, we describe the abstract message operations provided by the service.
In Section \ref{sec:model}, we end the document by presenting the data types used by the mentioned operations.
The concrete interface is specified in the HTTP/TLS/JSON IDD \cite{ca-authorizationToken-http-idd}. The MQTT binding is specified in the MQTT/TLS/JSON IDD \cite{ca-authorizationToken-mqtt-idd}.

\subsection{Significant Prior Art}

Successor of the Arrowhead 4.x token generation of the Authorization system (which issued JWTs during orchestration).

\subsection{How This Service Is Meant to Be Used}

Normally Dynamic\allowbreak{}Service\allowbreak{}Orchestration requests the token on behalf of the consumer during orchestration, when the provider interface has a token policy. A consumer may also call \texttt{generate} itself. The provider calls \texttt{verify} (simple tokens) or validates a JSON Web Token locally with the key from \texttt{get-\allowbreak{}public-\allowbreak{}key}. Providers that want encrypted self-contained tokens register an AES key with \texttt{register-\allowbreak{}encryption-\allowbreak{}key}.

\subsection{Important Delimitations and Dependencies}
\label{sec:delimitations}

A token is issued only if the consumer is authorized by the policies. Simple tokens are limited by time (\texttt{token.\allowbreak{}time.\allowbreak{}limit}) or number of uses (\texttt{simple.\allowbreak{}token.\allowbreak{}usage.\allowbreak{}limit}). The service does not transport tokens; the consumer sends them in its requests to the provider.

\newpage

\section{Service Interface}
\label{sec:operations}

This section describes the interfaces to the \texttt{authorizationToken} service.
In particular, each subsection names an abstract operation, an input type and an output type, in that order.
The input type is named inside parentheses, while the output type is preceded by a colon.
Input and output types are only denoted when accepted or returned, respectively, by the interface in question.

All abstract data types named in this section are defined in Section \ref{sec:model}.

\begin{figure}[ht!]
  \centering
  \begin{tikzpicture}[font=\small]
    \node[draw=ArrowheadBlue, very thick, rectangle split, rectangle split parts=2, align=center,
          rectangle split part fill={ArrowheadBlue!12,white}, text width=9cm, inner sep=6pt] (b)
      {\guillemotleft{}Service\guillemotright{}\\\textbf{authorizationToken}
        \nodepart{second} \texttt{generate}\\ \texttt{verify}\\ \texttt{get-\allowbreak{}public-\allowbreak{}key}\\ \texttt{register-\allowbreak{}encryption-\allowbreak{}key}\\ \texttt{unregister-\allowbreak{}encryption-\allowbreak{}key}};
  \end{tikzpicture}
  \caption{SysML-style block of the \texttt{authorizationToken} service and its operations.}
  \label{fig:SD}
\end{figure}

The following interface operations are available.

\osubsection{generate}{\mref[Authorization\allowbreak{}Token\allowbreak{}Generation\allowbreak{}Request]{AuthorizationTokenGenerationRequest}}{\mref[Authorization\allowbreak{}Token\allowbreak{}Generation\allowbreak{}Response]{AuthorizationTokenGenerationResponse}}

Generates a token of the requested variant for the requester (as consumer) to the provider's target and operation, if authorized.

\osubsection{verify}{\pref{Token}}{\mref[Authorization\allowbreak{}Token\allowbreak{}Verify\allowbreak{}Response]{AuthorizationTokenVerifyResponse}}

Verifies a token presented to the requester (as provider), counts a use of usage-limited tokens and returns its consumer, target and scope.

\osubsection{get-public-key}{}{\pref{String}}

Returns the public key (Base64 X.509 Subject\allowbreak{}Public\allowbreak{}Key\allowbreak{}Info) with which providers verify JSON Web Tokens.

\osubsection{register-encryption-key}{\mref[Authorization\allowbreak{}Encryption\allowbreak{}Key\allowbreak{}Registration\allowbreak{}Request]{AuthorizationEncryptionKeyRegistrationRequest}}{\pref{String}}

Registers or replaces the requester's AES key for encrypting its self-contained tokens. Returns the initialization vector (CBC) or an empty string; 201 for a new key.

\osubsection{unregister-encryption-key}{}{}

Removes the requester's encryption key. Returns 200 if removed, 204 if there was none.

\newpage

\section{Information Model}
\label{sec:model}

Here, all data objects that can be part of the \texttt{authorizationToken} service are listed in alphabetic order.
Note that each subsection, which describes one type of object, begins with the \textit{struct} keyword, which is used to denote a collection of named fields, each with its own data type.
As a complement to the explicitly defined types in this section, there is also a list of implicit primitive types in Section \ref{sec:model:primitives}, which are used to represent things like hashes and identifiers.

\msubsection{struct}{AuthorizationEncryptionKeyRegistrationRequest}

Registers the requester's key for encrypting its self-contained tokens.


\noindent\begin{tabularx}{\textwidth}{| >{\raggedright\arraybackslash}p{3.9cm} | >{\raggedright\arraybackslash\small}p{5.3cm} | X |} \hline
\rowcolor{gray!33} Field & Type & Description \\ \hline
key & \pref{String} & AES key of the provider (at least 16 bytes). \\ \hline
algorithm & \pref{String} & Encryption algorithm of the key, e.g. \texttt{AES/\allowbreak{}ECB/\allowbreak{}PKCS5Padding} or \texttt{AES/\allowbreak{}CBC/\allowbreak{}PKCS5Padding}. \\ \hline
\end{tabularx}



\msubsection{struct}{AuthorizationTokenGenerationRequest}

Request of the consumer for a token to one provider's service operation.


\noindent\begin{tabularx}{\textwidth}{| >{\raggedright\arraybackslash}p{3.9cm} | >{\raggedright\arraybackslash\small}p{5.3cm} | X |} \hline
\rowcolor{gray!33} Field & Type & Description \\ \hline
token\allowbreak{}Variant & \pref{String} & Requested token variant, as interface security policy (e.g. \texttt{TIME\_\allowbreak{}LIMITED\_\allowbreak{}TOKEN\_\allowbreak{}AUTH}, \texttt{RSA\_\allowbreak{}SHA512\_\allowbreak{}JSON\_\allowbreak{}WEB\_\allowbreak{}TOKEN\_\allowbreak{}AUTH}). \\ \hline
provider & \pref{Name} & Provider system the token is for. \\ \hline
target\allowbreak{}Type & \pref{String} & Target type: \texttt{SERVICE\_\allowbreak{}DEF} or \texttt{EVENT\_\allowbreak{}TYPE}. \\ \hline
target & \pref{Name} & Service definition (or event type). \\ \hline
scope & \pref{Name} & Operation the token is for. \\ \hline
\end{tabularx}



\msubsection{struct}{AuthorizationTokenGenerationResponse}

A generated token.


\noindent\begin{tabularx}{\textwidth}{| >{\raggedright\arraybackslash}p{3.9cm} | >{\raggedright\arraybackslash\small}p{5.3cm} | X |} \hline
\rowcolor{gray!33} Field & Type & Description \\ \hline
token\allowbreak{}Type & \pref{AuthorizationTokenType} & Token type: \texttt{TIME\_\allowbreak{}LIMITED\_\allowbreak{}TOKEN}, \texttt{USAGE\_\allowbreak{}LIMITED\_\allowbreak{}TOKEN}, \texttt{SELF\_\allowbreak{}CONTAINED\_\allowbreak{}TOKEN} or \texttt{TRANSLATION\_\allowbreak{}BRIDGE\_\allowbreak{}TOKEN}. \\ \hline
target\allowbreak{}Type & \pref{AuthorizationTargetType} & Target type: \texttt{SERVICE\_\allowbreak{}DEF} or \texttt{EVENT\_\allowbreak{}TYPE}. \\ \hline
token & \pref{Token} & The token (opaque string, Base64 or JSON Web Token). \\ \hline
usage\allowbreak{}Limit & \pref{Number} & Number of allowed uses (usage-limited tokens). \\ \hline
expires\allowbreak{}At & \pref{DateTime} & UTC expiry of time-limited tokens. \\ \hline
\end{tabularx}



\msubsection{struct}{AuthorizationTokenVerifyResponse}

Result of verifying a token presented to a provider.


\noindent\begin{tabularx}{\textwidth}{| >{\raggedright\arraybackslash}p{3.9cm} | >{\raggedright\arraybackslash\small}p{5.3cm} | X |} \hline
\rowcolor{gray!33} Field & Type & Description \\ \hline
verified & \pref{Boolean} & True when the token is valid. \\ \hline
consumer\allowbreak{}Cloud & \pref{Name} & Cloud identifier of the consumer system. \\ \hline
consumer & \pref{Name} & Name of the consumer system. \\ \hline
target\allowbreak{}Type & \pref{AuthorizationTargetType} & Target type: \texttt{SERVICE\_\allowbreak{}DEF} or \texttt{EVENT\_\allowbreak{}TYPE}. \\ \hline
target & \pref{Name} & Target of the policy or token: a service definition or event type. \\ \hline
scope & \pref{Name} & Operation the token or policy is limited to; empty for all operations. \\ \hline
\end{tabularx}



\subsection{Primitives}
\label{sec:model:primitives}

Types and structures mentioned throughout this document that are assumed to be available to implementations of this service.
The concrete interpretations of each of these types and structures must be provided by any IDD document claiming to implement this service.

\noindent\begin{tabularx}{\textwidth}{| p{4.6cm} | X |} \hline
\rowcolor{gray!33} Type & Description \\ \hline
\pdef{AuthorizationTargetType} & One of \texttt{SERVICE\_\allowbreak{}DEF}, \texttt{EVENT\_\allowbreak{}TYPE}. \\ \hline
\pdef{AuthorizationTokenType} & One of \texttt{TIME\_\allowbreak{}LIMITED\_\allowbreak{}TOKEN}, \texttt{USAGE\_\allowbreak{}LIMITED\_\allowbreak{}TOKEN}, \texttt{SELF\_\allowbreak{}CONTAINED\_\allowbreak{}TOKEN}, \texttt{TRANSLATION\_\allowbreak{}BRIDGE\_\allowbreak{}TOKEN}. \\ \hline
\pdef{Boolean} & One out of \texttt{true} or \texttt{false}. \\ \hline
\pdef{DateTime} & Pinpoints a specific moment in time (UTC). \\ \hline
\pdef{Name} & A string identifier that is intended to be both human and machine-readable. \\ \hline
\pdef{Number} & An integer number. \\ \hline
\pdef{String} & A chain of characters. \\ \hline
\pdef{Token} & An opaque string used as credential. \\ \hline
\end{tabularx}

\newpage

\bibliographystyle{IEEEtran}
\bibliography{bibliography}

\newpage

\section{Revision History}
\subsection{Amendments}

\noindent\begin{tabularx}{\textwidth}{| p{1cm} | p{3cm} | p{2cm} | X | p{4cm} |} \hline
\rowcolor{gray!33} No. & Date & Version & Subject of Amendments & Author \\ \hline

1 & 2026-10-22 & \arrowversion & Document created from the Arrowhead LaTeX templates for the 5.2.2 Java Spring implementation, including the local extensions of this workspace. & Jerker Delsing \\ \hline
\end{tabularx}

\subsection{Quality Assurance}

\noindent\begin{tabularx}{\textwidth}{| p{1cm} | p{3cm} | p{2cm} | X |} \hline
\rowcolor{gray!33} No. & Date & Version & Approved by \\ \hline

1 & & \arrowversion  &  \\ \hline

\end{tabularx}

\end{document}
